\subsection*{A finite-suspension test for the smallest reversible class}

\begin{claim}{D.2.7}{Proved}{Ordinary nullification after a bounded suspension}
\label{D27nullification}
Let $A$ be a connected pointed space whose reduced integral homology is
$p$-primary torsion. Suppose that, for some integer $r\geq1$, $A$ kills
every $r$-connected generalized Eilenberg--Mac Lane space (GEM) with
$p$-primary torsion homotopy groups. Then the aspherical class of the
smallest reversible fundamental localizer containing $A\to *$ is
\[
 \mathcal A_{\min\mathrm{Rev}(A)}
   =\{X\in\cS_*:P_A(\Sigma^rX)\simeq *\}.
 \tag{D.2.7a}
\]
Here $P_A$ is ordinary unstable $A$-nullification. In particular:
\begin{enumerate}
\item For $A=K(\mathbb F_p,n)$, $n\geq1$, one may take
      $r=\max\{1,n-1\}$.
\item For $A=K(C_{p^\infty},n)$, $n\geq1$, where
      $C_{p^\infty}=\mathbb Z[1/p]/\mathbb Z$, one may take $r=n$.
\item If $B\mathbb Z/p$ is $A$-acyclic, one may take $r=1$.
      Thus for every connected $V$ with $p$-primary reduced integral
      homology,
\[
 \mathcal A_{\min\mathrm{Rev}(B\mathbb Z/p\vee V)}
  =\{X:P_{B\mathbb Z/p\vee V}(\Sigma X)\simeq *\}.
 \tag{D.2.7b}
\]
\end{enumerate}
The statements impose no bound on cardinalities or torsion exponents
and use no $p$-completion.
\end{claim}

\begin{proof}
Put $\mathcal C_A=\{Y:P_AY\simeq *\}$ and
$\mathcal N_r=\{X:\Sigma^rX\in\mathcal C_A\}$.
Ordinary nullification exists for every space $A$
\citep[Theorem~1.1]{Bousfield1996ULP}. Its mapping-space universal
property shows that $\mathcal C_A$ is closed under weak equivalences,
retracts and pointed homotopy colimits. In a homotopy cofiber sequence,
membership of the first and second terms, or of the first and third
terms, implies membership of the remaining term: apply the derived
mapping-space functor into every $A$-null space. These same properties
hold for $\mathcal N_r$, since suspension preserves pointed homotopy
colimits and cofiber sequences.

Every $A$-acyclic space has $p$-primary reduced integral homology.
Indeed, $A\to *$ is an $H\mathbb Z[1/p]$-equivalence; every
$H\mathbb Z[1/p]$-local space is therefore $A$-null. The universal
property of $P_A$ implies that every $A$-acyclic space is
$H\mathbb Z[1/p]$-acyclic. Since $\mathbb Z[1/p]$ is flat, its
vanishing reduced homology means exactly that every element of each
reduced integral homology group is killed by a power of $p$.

We prove the essential desuspension implication. Suppose
$\Sigma X\in\mathcal N_r$, so that
$\Sigma^{r+1}X\in\mathcal C_A$. The preceding paragraph and the
suspension homology isomorphism imply that every
$\widetilde H_i(X;\mathbb Z)$ is $p$-primary torsion. In particular
$X$ is connected, because $\widetilde H_0(X;\mathbb Z)$ is free
abelian. Hence $\Sigma^rX$ is $r$-connected. By Dold--Thom,
$SP^\infty\Sigma^rX$ is an $r$-connected GEM with $p$-primary
torsion homotopy groups, and the hypothesis gives
$SP^\infty\Sigma^rX\in\mathcal C_A$.

The exact classical input is \citet[Theorem~3.3]{Bousfield1996ULP}:
for a connected space $U$ and an integer $k\geq1$,
\[
 \langle\Sigma U\rangle
  =\langle\Sigma^kU\vee SP^\infty\Sigma U\rangle.
\]
The brackets here denote ordinary nullity classes: equality says that
the two displayed spaces kill exactly the same spaces under ordinary
nullification. Apply this identity to the connected space
$U=\Sigma^{r-1}X$, with $k=2$. Both summands
$\Sigma^{r+1}X$ and $SP^\infty\Sigma^rX$ are $A$-acyclic, so their
wedge is $A$-acyclic. The identity gives
$\Sigma^rX\in\mathcal C_A$, proving $X\in\mathcal N_r$.

For the third cofiber implication, if $Y,C\in\mathcal N_r$ in a
cofiber sequence $X\to Y\to C$, use the next cofiber sequence
$Y\to C\to\Sigma X$ and the already proved first implication.
It gives $\Sigma X\in\mathcal N_r$, and desuspension gives
$X\in\mathcal N_r$. Thus A.3 constructs a reversible fundamental
localizer with aspherical class $\mathcal N_r$. It contains $A\to *$
because every positive suspension of $A$ is $A$-acyclic.

For minimality, let $W$ be any reversible fundamental localizer
containing $A\to *$. The ordinary localizer generated by this one map
is contained in $W$, so every space in $\mathcal C_A$ is
$W$-aspherical. If $X\in\mathcal N_r$, then
$\Sigma^rX\in\mathcal C_A$; reversible desuspension, applied $r$
times, puts $X$ in the aspherical class of $W$. The reversible cofiber
criterion identifies the localizers themselves. This proves
\textup{(D.2.7a)}.

For the first listed application, the opened statement of
\citet[Lemma~8.13]{EH5Bousfield1997} says that $K(G,n)$, for
$n\geq2$ and an arbitrary nonzero $p$-torsion abelian group $G$,
kills each $(n-1)$-connected generalized polyGEM with $p$-torsion
homotopy groups, subject to the additional condition
$\pi_nX/p\pi_nX=0$ only when $G/pG=0$. A GEM is a generalized
polyGEM. Taking $G=\mathbb F_p$ removes the additional condition,
and gives $r=n-1$ when $n\geq2$.

For $n=1$, $B\mathbb Z/p$ kills $K(\mathbb F_p,2)$ by the
simplicial bar construction: its levels are finite powers of
$B\mathbb Z/p$, nullification preserves finite products
\citep[proof of Theorem~1.7]{Bousfield1996ULP}, and its acyclic
class is closed under geometric realization. Applying Lemma~8.13 to
$K(\mathbb F_p,2)$ gives all simply connected $p$-torsion GEMs.
This proves $r=1$ and also proves the third listed application by
transitivity of ordinary acyclicity.

For $G=C_{p^\infty}$ and $n\geq2$, the extra condition in
Lemma~8.13 cannot be discarded: $G/pG=0$. It is automatically
satisfied on $n$-connected GEMs, giving $r=n$. To include $n=1$,
use the universal coefficient injection
\[
 \operatorname{Ext}_{\mathbb Z}(C_{p^\infty},\mathbb F_p)
  \lhook\joinrel\longrightarrow
 H^2(BC_{p^\infty};\mathbb F_p).
\]
The group on the left is nonzero: the extension
$0\to\mathbb F_p\to C_{p^\infty}\xrightarrow{p}C_{p^\infty}\to0$
does not split, since a direct summand of a divisible abelian group is
divisible. Thus $K(\mathbb F_p,2)$ is not
$BC_{p^\infty}$-null. By
\citet[Lemma~9.13]{EH5Bousfield1997}, an Eilenberg--Mac Lane
space $K(\mathbb F_p,m)$ is either acyclic or local for any specified
localization. It is therefore $BC_{p^\infty}$-acyclic. Lemma~8.13
again supplies all simply connected $p$-torsion GEMs, giving $r=1$.

Finally these Eilenberg--Mac Lane generators have $p$-primary
reduced integral homology. One can use the nilpotent localization
criterion away from $p$, or start with the finite cyclic groups and
use the bar construction, followed by filtered colimits for
$C_{p^\infty}$. This verifies the homology premise in every case.
\end{proof}

\begin{remark}
For the zero-finite-part example in D.2.6, the exact remaining unstable
test becomes
\[
 \Sigma^\infty X\in\loc(\Sigma^\infty K(\mathbb F_p,2)),
 \qquad P_{K(\mathbb F_p,2)}(\Sigma X)\not\simeq *.
\]
Likewise \textup{(D.2.7b)} turns the interval tests generated by
$B\mathbb Z/p\vee V$ into ordinary-nullification calculations.
Neither formula supplies the required separating space by itself;
D.2.4 and the strict-intermediate-element problem in D.3 retain
their stated status. The initial suspension is essential:
Theorem~3.3 makes no assertion that the unsuspended ordinary
$B\mathbb Z/p$-acyclic class is desuspension-closed.
\end{remark}

\begin{claim}{D.2.8}{Proved}{Finite primary space generators are $RT$-fixed}
Let $V$ be a nonzero finite connected pointed complex with $p$-primary
reduced integral homology, where ``nonzero'' means
$\Sigma^\infty V\not\simeq0$. Then
\[
 \min\mathrm{Rev}(V\to *)
   =R\bigl(\loc(\Sigma^\infty V)\bigr).
 \tag{D.2.8}
\]
This equality concerns all pointed source spaces, without a finiteness
hypothesis on them. Thus the finite primary generators themselves
give fixed localizers; the use of infinite-loop generators in D.2.2
is material.
\end{claim}

\begin{proof}
Put $W=\min\mathrm{Rev}(V\to *)$ and
$\mathcal L=\loc(\Sigma^\infty V)$. The inclusion
$W\subseteq R(\mathcal L)$ follows from the defining minimum.
It remains to prove the reverse inclusion on aspherical spaces.

Let $\Sigma^\infty X\in\mathcal L$. The spectrum $H\mathbb Z[1/p]$
annihilates $\Sigma^\infty V$, so it annihilates $\mathcal L$.
Consequently $X$ has $p$-primary reduced integral homology. Its
reduced $H_0$, being free abelian, is zero; hence $X$ is connected.
The space $Y=\Sigma X$ is simply connected and has $p$-primary
integral homology. The nilpotent homological localization criterion
therefore gives $Y\simeq Y_{(p)}$.

The nonzero finite $p$-primary spectrum $\Sigma^\infty V$ has some
finite type $n+1\geq1$. For each $0\leq i\leq n$, the telescope
$T(i)$ annihilates $\Sigma^\infty V$, and thus every spectrum in
$\mathcal L$. In particular $Y\to *$ is a $T(i)$-homology
equivalence for all these $i$.

We use the exact finite-generator theorem stated and proved in
\citet[Theorem~4.1]{FGHBarkanHeutsShi2026}, attributed there to
Bousfield's 1994 work. It says that, for an arbitrary map of spaces
$f$, being a $T(i)$-homology equivalence for every $i\leq n$ is
equivalent to
\[
 P_{\Sigma V}(\Sigma^k f_{(p)})
 \quad\hbox{being an equivalence for all sufficiently large }k.
\]
Apply it to $Y\to *$, using $Y\simeq Y_{(p)}$. We obtain
$P_{\Sigma V}(\Sigma^kY)\simeq *$ for some $k$.
The ordinary localizer generated by $\Sigma V\to *$ lies in $W$,
since $V$ and its suspension are $W$-aspherical. Thus
$\Sigma^{k+1}X$ is $W$-aspherical. Repeated reversible
desuspension puts $X$ in the same aspherical class. The reversible
cofiber criterion now gives the equality of localizers in
\textup{(D.2.8)}.

Only the acyclicity of finite type-$(n+1)$ spectra for $T(i)$,
$i\leq n$, and the displayed primary theorem are used. No
identification of telescopic and Morava-$K$ localizations is assumed.
\end{proof}

\begin{claim}{D.2.9}{Proved}{The connective stable obstruction and its limit}
Let $C\in\Ho(\Sp)_{\geq n}$, $n\geq1$, have $p$-primary integral
homology and $H_n(C;\mathbb F_p)\neq0$. Every
$E\in\loc(C)\cap\Ho(\Sp)_{\geq n}$ satisfies
\[
 P_C^{\Sp}E\simeq0,
 \tag{D.2.9a}
\]
where $P_C^{\Sp}$ is spectral nullification using the mapping
\emph{space} out of $C$, equivalently all its nonnegative
suspensions. Thus if $C=\Sigma^\infty A$, the infinite-loop space
$\Omega^\infty E$ is ordinarily $A$-acyclic.

In particular put $B=K(\mathbb F_p,2)$,
$\mathcal I_p=\loc(\Sigma^\infty B)$ and
$A_r=Q\Sigma^rB$, for $r\geq0$. Every
$E\in\mathcal I_p\cap\Ho(\Sp)_{\geq r+2}$ gives an ordinarily
$A_r$-acyclic infinite-loop space. The localizers
\[
 W_r=\min\mathrm{Rev}(A_r\to *)
\]
form a descending chain with common stable part $\mathcal I_p$,
which has zero finite part. The exact remaining test is
\[
 B\in\mathcal A_{W_r}
 \quad\Longleftrightarrow\quad
 P_{A_r}(\Sigma^{r+1}B)\simeq *.
 \tag{D.2.9b}
\]
Moreover, for $r\geq1$,
\[
 \begin{aligned}
 P_{A_r}A_{r-1}&\simeq K(\mathbb F_p,r+1),\\
 W_{r-1}&=W_r\vee
   \min\mathrm{Rev}(K(\mathbb F_p,r+1)\to *).
 \end{aligned}
 \tag{D.2.9c}
\]
The connective assertion does not decide this test, and neither a
strict inclusion in this chain nor its collapse is asserted.
\end{claim}

\begin{proof}
Write $C_i=\Sigma^{-i}C$ and $Y_i=P_{C_i}^{\Sp}E$, for
$i\geq0$. Their natural maps point from $Y_i$ to $Y_{i+1}$.
The precise input is \citet[Theorem~2.6]{Bousfield1996ULP}:
because $C$ satisfies the $n$-supported $p$-torsion condition,
\[
 \fib(Y_i\longrightarrow Y_{i+1})
   \simeq\Sigma^{n-i-1}HG_i
 \tag{D.2.9d}
\]
for some $p$-primary torsion abelian group $G_i$. Hence, for
$j>i$, the fiber of $Y_i\to Y_j$ lies in
$\Ho(\Sp)_{\leq n-i-1}$, by extension closure. Set
$Y_\infty=\colim_iY_i$. Filtered colimits in spectra are exact
and preserve these homotopy-group bounds, so
\[
 \fib(Y_i\to Y_\infty)\in\Ho(\Sp)_{\leq n-i-1}.
 \tag{D.2.9e}
\]

Fix $i$, and use the map $Y_{i+1}\to Y_\infty$. Its fiber lies
in $\Ho(\Sp)_{\leq n-i-2}$, and its suspended fiber lies in
$\Ho(\Sp)_{\leq n-i-1}$. On the other hand
$C_i\in\Ho(\Sp)_{\geq n-i}$. The $t$-structure orthogonality
therefore shows that
\[
 \Map_{\Sp}(C_i,Y_{i+1})
   \xrightarrow{\simeq}\Map_{\Sp}(C_i,Y_\infty).
\]
The first mapping space is contractible because $Y_{i+1}$ is
$C_{i+1}$-null and $C_i=\Sigma C_{i+1}$. Thus $Y_\infty$ is
null for every positive and negative shift of $C$. The shift from
$Y_i$ to $Y_{i+1}$ in this paragraph controls components as well
as positive homotopy groups; a vanishing fiber alone would not
have sufficed.

The cone-attachment construction of spectral nullification puts
each $\cofib(E\to Y_i)$ in $\loc(C)$. Since $E\in\loc(C)$,
also $Y_i$ and $Y_\infty$ belong to $\loc(C)$. A member of
$\loc(C)$ that is right orthogonal to all shifts of $C$ is zero:
its mapping spectrum out of every member of $\loc(C)$ vanishes,
including its own identity. Hence $Y_\infty=0$.
Equation~\textup{(D.2.9e)} with $i=0$ now gives
$Y_0\in\Ho(\Sp)_{\leq n-1}$. But nullification by
$n$-connective $C$ preserves $n$-connectivity, since attaching
cones on nonnegative suspensions of $C$ stays in
$\Ho(\Sp)_{\geq n}$. Thus $Y_0=0$, proving
\textup{(D.2.9a)}. This argument neither assumes $C$ compact
nor interchanges its mapping-space functor with a filtered
colimit.

The passage to infinite-loop spaces is exactly
\citet[Theorem~2.10 and Proposition~3.1]{Bousfield1996ULP}:
\[
 P_A\Omega^\infty E
   \simeq\Omega^\infty P_{\Sigma^\infty A}^{\Sp}E.
\]
For $A_r$, Snaith splitting identifies
$\loc(\Sigma^\infty A_r)=\mathcal I_p$: its linear summand is
$\Sigma^r\Sigma^\infty B$, and all higher extended powers
belong to its localizing tensor ideal. The first homology group
of $A_r$ is $H_{r+2}(A_r;\mathbb Z)=\mathbb F_p$, and
$\Sigma^\infty A_r$ is $(r+2)$-connective. The already proved
assertion therefore applies with $n=r+2$.

The connected infinite-loop space $A_{r+1}$ is the delooping of
$A_r$. Its simplicial bar construction, preservation of finite
products by ordinary nullification, and closure of its acyclic
class under realization show that $A_r$ kills $A_{r+1}$.
Consequently $W_{r+1}\subseteq W_r$. Also
$W_0=\min\mathrm{Rev}(B\to *)$: the infinite-loop structure
on $B$ makes it a retract of $QB$, while every space kills its
free infinite-loop space by Proposition~3.1 of the same source.
Snaith splitting gives $T(W_r)=\mathcal I_p$ for every $r$,
and D.2.6 gives its zero finite part.

We compute the adjacent ordinary nullification explicitly. First
$H\mathbb F_p\in\mathcal I_p$: smash its generator with
$H\mathbb F_p$ and split the resulting nonzero
$H\mathbb F_p$-module into a wedge of shifts of
$H\mathbb F_p$. Its degree-two homology supplies a nonzero
summand, and $\mathcal I_p$ is closed under retracts and shifts.
For $r\geq1$, set $E=\Sigma^{r-1}\Sigma^\infty B$. Its lower
Postnikov section is
\[
 \tau_{\leq r+1}E\simeq\Sigma^{r+1}H\mathbb F_p.
\]
Both $E$ and this section belong to $\mathcal I_p$, so their
fiber $\tau_{\geq r+2}E$ does also. Apply
\textup{(D.2.9a)} with $C=\Sigma^\infty A_r$ and $n=r+2$ to
kill that fiber. The displayed lower section is already
$C$-null by $t$-structure orthogonality. Spectral nullification
preserves this cofiber sequence because its third term is
$C$-null \citep[Corollary~2.3]{Bousfield1996ULP}. Hence
\[
 P_{\Sigma^\infty A_r}^{\Sp}E
   \simeq\Sigma^{r+1}H\mathbb F_p.
\]
Theorem~2.10 then gives the first equality in
\textup{(D.2.9c)}, including its natural nullification unit
$A_{r-1}\to K(\mathbb F_p,r+1)$. This unit is a
$W_r$-equivalence. Consequently adjoining the collapse of its
target to $W_r$ is equivalent to adjoining the collapse of
$A_{r-1}$. Since $W_r\subseteq W_{r-1}$, the second equality
in \textup{(D.2.9c)} follows. Thus an adjacent strictness question
is reduced to the membership of one specified
Eilenberg--Mac Lane space, rather than settled by a connectivity
comparison.

Finally $A_r$ kills its infinite symmetric product, which has
$K(\mathbb F_p,r+2)$ as a homotopy retract. Lemma~8.13 of
\citet{EH5Bousfield1997} shows that it kills every
$(r+1)$-connected $p$-torsion GEM. Apply D.2.7 with the bound
$r+1$ to obtain \textup{(D.2.9b)}. That target is a suspension
of $B$; the spectral argument does not supply an infinite-loop
description of it. This is the exact unproved step separating
the reduction from a zero-finite-part counterexample.
\end{proof}
